\section{Related work}
\label{sec:related}

We review four streams: the lot-sizing and scheduling taxonomy in
which the PSP is embedded (Section~\ref{sec:rw-taxonomy}); applications
featuring process selection and co-production in fixed proportions
(Section~\ref{sec:rw-coproduction}); MIP-based matheuristics of the
relax-and-fix and fix-and-optimize families
(Section~\ref{sec:rw-matheuristics}); and the solver-progress and
experimental-methodology literature that underpins our protocol
(Section~\ref{sec:rw-methodology}).

\subsection{Lot sizing and scheduling: taxonomy and reviews}
\label{sec:rw-taxonomy}

Dynamic lot-sizing models that integrate sequencing decisions are
classically organized by their time structure. Small-bucket models
allow at most one (or a few) setups per period: the discrete
lot-sizing and scheduling problem \citep[DLSP;][]{fleischmann1990dlsp}
imposes all-or-nothing production, the continuous setup lot-sizing
problem \citep[CSLP;][]{karmarkar1985deterministic} relaxes full
capacity usage, and the proportional lot-sizing and scheduling problem
\citep[PLSP;][]{drexl1995proportional} allows the setup state to
change once within a period. The general lot-sizing and scheduling
problem \citep[GLSP;][]{fleischmann1997general} subsumes these
structures through micro-periods, and the capacitated lot-sizing
problem with sequence-dependent setups \citep{haase1996clsd} extends
the big-bucket tradition. Comprehensive surveys trace the evolution of
this family: \citet{drexl1997lotsizing} and
\citet{karimi2003capacitated} for the foundations,
\citet{jans2008modeling} for industrial extensions,
\citet{buschkuhl2010dynamic} for solution approaches,
\citet{glock2014tertiary} at the tertiary level, and, most relevant to
our positioning, the classification of simultaneous lot-sizing and
scheduling models by \citet{copil2017simultaneous}, recently extended
to secondary resources by \citet{worbelauer2019secondary}. Within this
taxonomy the PSP is a single-stage, single-machine, big-bucket model
with at most one ``product'' -- in fact, one \emph{process} -- per
period; its distinctive trait, absent from the standard classes, is
that the assignment decision releases a technologically fixed
\emph{vector} of products rather than a single item. Demand is met
against cumulative order-book quantities with shortage minimization,
which relates the model to service-level-oriented variants rather than
to the cost-minimization mainstream \citep{pochet2006production}.

\subsection{Process selection and co-production in fixed proportions}
\label{sec:rw-coproduction}

The PSP was introduced in the electrofused grains context by
\citet{luche2009combining}, who combined process selection with
single-stage lot sizing, proposed a constructive heuristic, and
reported that instances of practical size defeated the MIP technology
of the time. A structurally similar coupling of process configuration,
lot sizing, and scheduling arises in the molded pulp industry
\citep{martinez2018coupled}, where each mold configuration also
determines a mix of items, and in small foundries, where each furnace
alloy enables a set of castable parts \citep{araujo2008foundry}.
Related recipe- or process-based settings include animal nutrition
plants \citep{toso2009lot,clark2010production}, chemical plants in
which one item may be produced by alternative processes
\citep{villasboas2021chemical}, and integrated soft-drink production
where tank lots feed bottling lines
\citep{ferreira2009solution,toledo2015sitlsp}. Co-production in fixed
(or stochastic) proportions is also the defining feature of
semiconductor planning, where each wafer lot ``bins'' into multiple
performance grades \citep{leachman2002semiconductor,ng2007robust},
and connects to multi-period cutting-stock models in which a cutting
pattern plays the role of a process \citep{silva2023formulations}.
Across this literature, exact approaches are typically abandoned at
realistic scales in favor of tailored heuristics; none of these works,
however, revisits the abandoned monoliths with current solver
technology, and public benchmarks for fixed-proportion co-production
are, to the best of our knowledge, unavailable.

\subsection{Relax-and-fix and fix-and-optimize matheuristics}
\label{sec:rw-matheuristics}

Relax-and-fix constructs solutions by partitioning the integer
variables -- most naturally along the time axis -- and solving a
sequence of MIPs in which one window is kept integral, earlier windows
are fixed, and later windows are relaxed
\citep{dillenberger1994practical,pochet2006production}.
Fix-and-optimize (also called exchange or MIP-based improvement)
iteratively re-optimizes subsets of integer variables while fixing the
remainder at incumbent values; it was popularized for multi-level
capacitated lot sizing by \citet{sahling2009solving} and
\citet{helber2010fix}. The combination of the two -- \RF{} for
construction, \FO{} for improvement -- was systematized by
\citet{toledo2015relax} on multi-level lot-sizing and glass-container
problems, and industrial variants have since accumulated: integrated
pulp and paper planning \citep{santos2012integrated}, later improved
by variable neighborhood search \citep{figueira2013hybrid} and by
structure-exploiting matheuristics \citep{furlan2024matheuristic};
beverage plants with temporal cleanings \citep{toscano2020formulation};
and perishable-product lines, where relax-and-fix also supplies
competitive dual bounds \citep{jors2021perishable}. Iterative
MIP-based neighborhood searches for lot sizing and scheduling were
explored by \citet{james2011single}, and the genre is surveyed under
the ``matheuristics'' label by \citet{fischetti2018matheuristics}.
Our \RFFO{} follows the time-window template of
\citet{toledo2015relax}; the differentiating elements are the
co-production structure (windows re-optimize the process schedule,
with all product balances linked through cumulative demand), a fully
specified budget-control policy whose implementation details we report
for reproducibility, and the \RFMIP{} variant in which the \RF{}
solution warm-starts the \emph{monolithic} model -- a hybrid whose
value, as our experiments show, depends on whether the incumbent it
buys is worth the branch-and-bound time it costs.

\subsection{Solver progress and experimental methodology}
\label{sec:rw-methodology}

Machine-independent measurements document speedups of several orders
of magnitude in MIP solving between the early 2000s and 2020, driven
by presolve, cutting planes, primal heuristics, and parallel search
\citep{koch2022progress,achterberg2013analyzing,achterberg2020presolve,bixby2012brief}.
This progress reframes older intractability findings: conclusions
drawn against CPLEX~11--12 -- as in the original PSP studies -- need
not hold against CPLEX~22, and revisiting them yields calibrated
evidence of practical value. On the methodology side, our experimental
protocol adopts community standards for heuristic--exact comparisons:
performance profiles \citep{dolan2002benchmarking}, time-to-target
plots for randomized methods \citep{aiex2007ttt}, and public release
of instances and solutions in the spirit of MIPLIB
\citep{gleixner2021miplib}. The metaheuristic baseline itself follows
established components: GRASP \citep{feo1995greedy} with reactive
parameter control \citep{prais2000reactive}, iterated local search
\citep{lourenco2010iterated}, and path relinking
\citep{laguna1999grasp,resende2016optimization}. Finally, the
managerial framing of short re-planning budgets draws on the
rolling-horizon tradition initiated by \citet{baker1977experimental}.

\medskip
\noindent\textbf{Gap.} In summary: (i) the PSP with co-production in
fixed proportions has not been revisited since 2011, despite solver
progress that plausibly changes its practical status; (ii) matheuristic
studies seldom report primal--primal comparisons against a modern
solver under strictly equal wall-clock and thread budgets, leaving the
``method of choice'' question anecdotal; and (iii) no public benchmark
exists for this problem class. This paper addresses all three gaps,
with hypotheses registered before experimentation.
